\section{Certified cuts and their replay}\label{sec:certification}

A separation routine has two parts. One part chooses a direction; the other
proves a lower bound for that direction. Only the second affects validity.
This section states what must be checked for a cut in the original variables
to be valid as the solver stores it, and how a record of the check is
replayed.

\subsection{The stored row}\label{sec:exported-row}

Floating-point LP solvers store binary64 coefficients. Every finite binary64
number is a dyadic rational, so the row that the solver uses is a rational
inequality, and validity can be stated for that row. For a cut of the
form~\eqref{eq:agg-cut}, four things must hold.
\begin{enumerate}[label=(C\arabic*)]
\item The direction $(a,\lambda)$ is taken as the exact rational value of its
binary64 representation, and a lower bound $L\le\varphi(a,\lambda)$ over the
whole domain $D$ is certified.
\item The cut is assembled exactly: $c=S\T a-A\T\lambda$ and $r=L-\lambda\T b$,
with repeated occurrences of a variable combined.
\item The rational row $c\T v\ge r$ is rounded to binary64 with a correction
that keeps it valid (Proposition~\ref{prop:export}).
\item The row that the solver stores is the exported row.
\end{enumerate}
By Proposition~\ref{prop:aggregate}, (C1) and (C2) give a valid rational row,
and (C3) and (C4) carry its validity to the stored row. The direction itself
may come from an LP over samples, from a heuristic or from a previous
iteration; none of these needs to be correct. The principle is established:
a numerically chosen affine bound needs no verification of how it was found,
only a rigorous final shift \citep[Section~5]{GarloffJanssonSmith2003JCAM},
\citep{GarloffSmith2008}; cut parameters may be chosen in floating point and
the derivation then repeated rigorously
\citep[p.~294]{NeumaierShcherbina2004}; and safe linear estimators are
defined in terms of the machine coefficients themselves
\citep{BorradaileVanHentenryck2005}. The hazard at the last step is known to
SCIP's developers. Outside its exact mode, SCIP~10 rounds a row coefficient
that is integral within its tolerance to that integer and drops coefficients
below its epsilon, without adjusting the sides (functions \code{rowAddCoef},
\code{rowChgCoefPos} and \code{rowMerge} in \code{src/scip/lp.c} of
\citealp{SCIPsource10}). Its nonlinear constraint handler therefore relaxes its
own cuts with variable bounds before storing them, in floating point, and
rounds without correction when a needed bound is infinite
\citep[Section~4.2.10]{SCIP8report}; cuts added by another separator receive
no such treatment. Two points are specific to support cuts: the bound must be
certified for the final direction on the whole block domain, which is a
global nonconvex minimization, and the row that SCIP stores must be compared
with the certified row.

\begin{proposition}[Safe export]\label{prop:export}
Let $c\T v\ge r$ be valid, let $\hat c$ be binary64 coefficients and
$\Delta=\hat c-c$, and let $\ell_j\le v_j\le u_j$ be valid bounds, where
$\ell_j$ is finite whenever $\Delta_j>0$ and $u_j$ is finite whenever
$\Delta_j<0$. Put
$E=\sum_{j:\Delta_j>0}\Delta_j\ell_j+\sum_{j:\Delta_j<0}\Delta_ju_j$. Then
$\hat c\T v\ge\hat r$ is valid for every $\hat r\le r+E$.
\end{proposition}

\begin{proof}
$\hat c\T v=c\T v+\Delta\T v\ge r+E$.
\end{proof}

This is the bound-corrected rounding of safe LP bounds and safe cuts
\citep{NeumaierShcherbina2004,CookEtAl2009SafeCuts,EiflerGleixner2024},
applied to the eliminated row. A coefficient that is exactly representable
needs no bound, even on an unbounded variable such as the objective variable.
The correction must be applied after every transformation that changes the
row, including merging, scaling and the removal of tiny coefficients; $\hat r$
is rounded downward and compared with $r+E$ in exact arithmetic. If a needed
bound is infinite or the correction destroys the violation, the row is
discarded; a fixed safety shift does not replace the calculation
(Section~\ref{sec:validity} gives an example). The bounds
must be global implications of the model (Section~\ref{sec:faithful}). Our
exporter is more conservative than the proposition: it requires both bounds
of a variable to be finite whenever the coefficient of that variable changes.
For support inequalities in lifted coordinates $(x,w)$, as in our first
implementation, bounds on $w$ need not be available, and the direct form is
used instead: if $L\le\min_{x\in D}\{\hat a\T x+\hat b\T F(x)\}$ for the exact
values of the stored coefficients $(\hat a,\hat b)$ and $\hat r\le L$, then
$\hat a\T x+\hat b\T w\ge\hat r$ is valid for $H_D(F)$ by
Proposition~\ref{prop:support}.

\subsection{Lower bounds on the whole domain}\label{sec:bounds}

For quadratic blocks, (C1) is computed exactly by the oracles of
Section~\ref{sec:quadratic}. For polynomial \emph{remainders} (the
nonlinear part of a row, Section~\ref{sec:blocks}) in one or two
variables we use Bernstein enclosures on a subdivision of the block box,
computed in exact rational arithmetic
\citep{GarloffJanssonSmith2003,GarloffSmith2008,MunozNarkawicz2013}; for
univariate elementary expressions, outward-rounded ball arithmetic
\citep{Johansson2017arb} on interval cells, strengthened by a chord
correction from a verified bound on the second derivative
(Appendix~\ref{app:bounds}). A certificate depends on two hypotheses. The
cells must cover the box exactly, starting from the trusted box and not from
a list supplied by the generator, and a bound proved on a smaller domain is
never reused on a larger one. Domain requirements, such as a positive
argument of $\log$ or a nonzero denominator, are proved on every cell; they
are collected from the raw source expression before any symbolic
simplification (Section~\ref{sec:faithful}), and every original feature is
enclosed on the whole cell before features are combined, so that
cancellation between features cannot hide a requirement. The chord
correction is used only on cells where the expression is twice continuously
differentiable and a finite bound on its second derivative is verified. A
bound computation returns one of four outcomes: a certified bound with its
partition; \emph{empty}, when every cell is excluded by a domain row, so that
$D=\emptyset$; \emph{unsupported}, for an operation outside the grammar; or
\emph{incomplete}, when a work limit is reached. Only the first produces a
cut. An incomplete computation says nothing about the point being separated;
in particular, it is not evidence that the point lies in the hull.

\subsection{Records and replay}\label{sec:replay}

Each cut is accepted through the same sequence: certify the support bound for
the binary64 direction (C1); eliminate the nonlinear rows exactly and merge
coefficients (C2); apply the correction of Proposition~\ref{prop:export} (C3);
and read back the row that SCIP has created from the cut, with its columns,
coefficients, constant, sides and scope, before it is passed to SCIP's
separation storage, and compare it with the certified row after mapping
original to transformed variables (C4). Any difference, such as a coefficient
rounded by SCIP or a variable aggregated or removed by presolve, rejects the
cut. Accepted rows are passed to SCIP as global, removable cuts with
\code{forcecut} set. SCIP's later handling of the row, such as converting a
cut in a single variable into a bound change, is outside this check.

The record of a cut contains the serialized source model, the signed source
rows and multipliers, the proven bounds with their derivations, the support
domain, the direction, the support certificate (the complete candidate
enumeration, the star partition or the subdivision with its bounds), the
eliminated row, the rounding correction and the stored SCIP row. Replay runs
in a fresh process. It verifies the hashes of the archived sources and model,
re-executes the model builder (which constructs, but does not solve, a SCIP
model through PySCIPOpt) and the bound derivations, rebuilds every signed
source side from the archived model with its own affine splitter, repeats the
exact support calculation for the recorded binary64 direction on the rebuilt
features and domain, recomputes elimination and rounding, and compares the
result with the stored row as recorded by the separator; the comparison with
SCIP's row itself is made only when the cut is generated. For one recorded cut in each part and cut mode,
fourteen single-field corruptions check that replay rejects altered inputs:
the model hash, the support normal, the right-hand side and identity of a
source side, a feature, the box, a row right-hand side, the scope, the
rounding correction, the support identity, a stored coefficient, a stored
side, the column mapping and the bound proof.

Replay re-executes the generator's exact code: the model builder, the bound
replay, the support kernels and the elimination and export routine. It does
not run the direction search, the LP solver or SCIP's solving process, and it
has its own affine splitter. It is therefore a re-execution in a fresh
process, not an independent checker. The trusted base consists of Python's
rational arithmetic; SymPy's expression construction (including its automatic
canonicalization), expansion, polynomial conversion and differentiation; Arb
through python-flint; the OSiL reader, both its parsing and its conversion of
decimal tokens to binary64; PySCIPOpt's translation of the checked expression
objects and row sides into SCIP constraints, and its accessors for stored
rows; and the code shared by generation and replay. None of these is formally
verified. The certificates say nothing about SCIP's presolve, its LP bounds,
its branching or its final optimality claim; SCIP's exact mode, which solves
MILPs without tolerances and can write VIPR certificates of complete solves
for independent checking \citep{CheungGleixnerSteffy2017,EiflerGleixner2023},
is limited to MILP \citep{SCIP10}.
